\documentclass[11pt,letterpaper]{article}
\usepackage[spanish,es-nodecimaldot]{babel}
\usepackage[utf8]{inputenc}
\usepackage[T1]{fontenc}
\usepackage{lmodern,microtype}
\usepackage[margin=25mm,inner=30mm,top=30mm]{geometry}
\usepackage{setspace}
\setstretch{1.3}
\setlength{\parindent}{0pt}
\setlength{\parskip}{5pt}
\usepackage{xcolor,graphicx,float,booktabs,caption}
\usepackage{fancyhdr}
\usepackage[colorlinks=true,linkcolor=black,urlcolor=blue]{hyperref}
\usepackage{xurl}
\usepackage{listings}
\usepackage[section]{placeins}
\definecolor{azulusm}{HTML}{174B79}
\definecolor{sqlbg}{HTML}{F4F6F9}
\lstdefinelanguage{MySQLLab}{
 language=SQL,
 morekeywords={AUTO_INCREMENT,DECIMAL,ENUM,SHOW,SET,SQL_SAFE_UPDATES,COALESCE},
 sensitive=false
}
\lstdefinestyle{sqlusm}{
 language=MySQLLab,basicstyle=\ttfamily\footnotesize,
 keywordstyle=\color{azulusm}\bfseries,
 commentstyle=\color{gray},stringstyle=\color{teal!65!black},
 breaklines=true,breakatwhitespace=false,
 backgroundcolor=\color{sqlbg},frame=single,rulecolor=\color{black!25},
 columns=fullflexible,keepspaces=true,showstringspaces=false,
 xleftmargin=3pt,xrightmargin=3pt,
 literate={á}{{\'a}}1 {é}{{\'e}}1 {í}{{\'i}}1 {ó}{{\'o}}1 {ú}{{\'u}}1
 {Á}{{\'A}}1 {É}{{\'E}}1 {Í}{{\'I}}1 {Ó}{{\'O}}1 {Ú}{{\'U}}1
 {ñ}{{\~n}}1 {Ñ}{{\~N}}1 {ü}{{\"u}}1
}
\lstset{style=sqlusm}
\newcommand{\theCampus}{Campus Vitacura}
\pagestyle{fancy}
\fancyhf{}
\fancyhead[L]{\IfFileExists{Logo_UTFSM.png}{\includegraphics[height=0.9cm]{Logo_UTFSM.png}}{\small\bfseries Universidad Técnica Federico Santa María}}
\fancyfoot[L]{\scriptsize Laboratorio N.º 5}
\fancyfoot[C]{\thepage}
\fancyfoot[R]{\scriptsize ICN292 -- USM}
\setlength{\headheight}{36pt}
\renewcommand{\headrulewidth}{0.4pt}
\renewcommand{\footrulewidth}{0.4pt}
\captionsetup{font=small,labelfont=bf}
\newcommand{\evidencia}[3][0.95]{
  \begin{figure}[H]\centering
  \includegraphics[width=#1\linewidth,height=0.27\textheight,keepaspectratio]{capturas/#2.png}
  \caption{#3}\end{figure}
}
\begin{document}
\begin{titlepage}
\thispagestyle{empty}
\begin{center}
\vspace*{1cm}
\IfFileExists{header utfsm.png}{\includegraphics[width=0.60\textwidth]{header utfsm.png}}{
{\Large\bfseries UNIVERSIDAD TÉCNICA FEDERICO SANTA MARÍA}\\[0.25cm]
{\large Departamento de Industrias — Campus Vitacura}}
\vspace{1.6cm}
\par\noindent\rule{\linewidth}{0.8pt}\par\vspace{0.2cm}
{\LARGE\bfseries Informe de Laboratorio N.º 5}\\[0.35cm]
{\Large Bases de Datos y SQL — Andes Bike}\\[0.35cm]
{\large Sistemas de Información para la Gestión (ICN292)}
\vspace{0.2cm}\par\noindent\rule{\linewidth}{0.8pt}\par
\vspace{1.5cm}
{\Large\bfseries Sergio Martínez}\\[0.5cm]
{\large Semestre 2026/02}\\[0.35cm]
{\large Departamento de Industrias}\\[0.5cm]
\vspace{0.7cm}
\textbf{Repositorio público GitHub (verificar publicación):}\\[0.2cm]
\url{https://github.com/Sergi9787/ICN292-Lab5-Martinez-Sergio}\\
\vfill
{\large Noviembre 2026}\\[0.35cm]
{\large\theCampus}
\end{center}
\end{titlepage}
\section*{Entregable 2: Consultas SQL}
\subsection*{Consulta 1. Vendedores y cantidad de ventas}
\begin{lstlisting}
SELECT vd.nombre AS vendedor,
       COUNT(v.id_venta) AS cantidad_ventas
FROM vendedores vd
LEFT JOIN ventas v ON v.id_vendedor = vd.id_vendedor
GROUP BY vd.id_vendedor, vd.nombre
ORDER BY cantidad_ventas DESC, vd.nombre ASC;
\end{lstlisting}
\evidencia[0.75]{c01}{Cantidad de ventas por vendedor, incluidos los vendedores sin ventas.}
\FloatBarrier

\subsection*{Consulta 2. Ventas y unidades vendidas por tienda}
\begin{lstlisting}
SELECT t.nombre AS tienda,
       COUNT(v.id_venta) AS cantidad_ventas,
       COALESCE(SUM(v.unidades), 0) AS unidades_vendidas
FROM tiendas t
LEFT JOIN vendedores vd ON vd.id_tienda = t.id_tienda
LEFT JOIN ventas v ON v.id_vendedor = vd.id_vendedor
GROUP BY t.id_tienda, t.nombre
ORDER BY unidades_vendidas DESC;
\end{lstlisting}
\evidencia[0.8]{c02}{Ventas y unidades de cada tienda, incluida Andes La Serena.}
\FloatBarrier

\subsection*{Consulta 3. Tiendas sin vendedores asignados}
\begin{lstlisting}
SELECT t.nombre AS tienda
FROM tiendas t
LEFT JOIN vendedores vd ON vd.id_tienda = t.id_tienda
WHERE vd.id_vendedor IS NULL;
\end{lstlisting}
\evidencia[0.55]{c03}{Tienda sin vendedores asignados.}
\FloatBarrier

\subsection*{Consulta 4. Modelos sin ventas en junio de 2025}
\begin{lstlisting}
-- El filtro temporal está en ON para conservar los modelos sin coincidencias.
SELECT b.modelo
FROM bicicletas b
LEFT JOIN ventas v
  ON v.id_bicicleta = b.id_bicicleta
 AND v.fecha >= '2025-06-01'
 AND v.fecha < '2025-07-01'
WHERE v.id_venta IS NULL
ORDER BY b.modelo;
\end{lstlisting}
\textbf{Explicación solicitada:} El filtro de junio de 2025 se colocó en la condición \texttt{ON} del \texttt{LEFT JOIN}, para mantener los modelos sin ventas en ese mes. Si se ubicara en \texttt{WHERE}, se excluirían las filas sin coincidencia y no se identificarían correctamente los modelos no vendidos.

\evidencia[0.6]{c04}{Modelos sin ventas registradas durante junio de 2025.}
\FloatBarrier

\subsection*{Consulta 5. Vendedores que no registraron ventas}
\begin{lstlisting}
SELECT vd.nombre AS vendedor,
       t.nombre AS tienda
FROM vendedores vd
LEFT JOIN ventas v ON v.id_vendedor = vd.id_vendedor
LEFT JOIN tiendas t ON t.id_tienda = vd.id_tienda
WHERE v.id_venta IS NULL
ORDER BY vd.nombre;
\end{lstlisting}
\evidencia[0.7]{c05}{Vendedores sin ventas y nombre de su tienda, cuando existe.}
\FloatBarrier

\subsection*{Consulta 6. Pares de vendedores de una misma tienda}
\begin{lstlisting}
-- El más antiguo primero; ante empate, menor id_vendedor primero.
SELECT a.nombre AS vendedor_1,
       a.fecha_ingreso AS ingreso_1,
       b.nombre AS vendedor_2,
       b.fecha_ingreso AS ingreso_2,
       t.nombre AS tienda
FROM vendedores a
JOIN vendedores b
  ON a.id_tienda = b.id_tienda
 AND (a.fecha_ingreso < b.fecha_ingreso
      OR (a.fecha_ingreso = b.fecha_ingreso AND a.id_vendedor < b.id_vendedor))
JOIN tiendas t ON t.id_tienda = a.id_tienda
ORDER BY t.nombre, a.fecha_ingreso, b.fecha_ingreso, a.nombre, b.nombre;
\end{lstlisting}
\evidencia[0.94]{c06}{Pares únicos de vendedores de una misma tienda, con sus fechas de ingreso.}
\FloatBarrier

\subsection*{Consulta 7. Monto vendido de las categorías Montaña y Ruta}
\begin{lstlisting}
SELECT t.nombre AS tienda,
       SUM(v.unidades * b.precio) AS monto_total
FROM tiendas t
JOIN vendedores vd ON vd.id_tienda = t.id_tienda
JOIN ventas v ON v.id_vendedor = vd.id_vendedor
JOIN bicicletas b ON b.id_bicicleta = v.id_bicicleta
WHERE b.categoria IN ('Montaña', 'Ruta')
GROUP BY t.id_tienda, t.nombre
HAVING SUM(v.unidades * b.precio) > 2000000
ORDER BY monto_total DESC;
\end{lstlisting}
\evidencia[0.75]{c07}{Tiendas cuyo monto de ventas de Montaña y Ruta supera los \$2.000.000.}
\FloatBarrier

\subsection*{Consulta 8. Bicicletas con precio superior al promedio}
\begin{lstlisting}
SELECT b.modelo, b.precio
FROM bicicletas b
WHERE b.precio > (SELECT AVG(precio) FROM bicicletas)
ORDER BY b.precio DESC;
\end{lstlisting}
\textbf{Explicación solicitada:} El precio promedio del catálogo es \$657.490, obtenido mediante \texttt{AVG(precio)} sobre los ocho registros de la tabla \texttt{bicicletas}. La subconsulta utiliza ese promedio para filtrar los modelos más caros.

\evidencia[0.65]{c08}{Modelos de bicicletas con precio superior al promedio del catálogo.}
\begin{lstlisting}
SELECT AVG(precio) AS precio_promedio_catalogo
FROM bicicletas;
\end{lstlisting}
\evidencia[0.60]{promedio_catalogo}{Comprobación auxiliar del precio promedio del catálogo: \$657.490.}
\FloatBarrier

\subsection*{Consulta 9. Vendedores con promedio de unidades superior al general}
\begin{lstlisting}
SELECT vd.nombre AS vendedor,
       AVG(v.unidades) AS promedio_unidades
FROM vendedores vd
JOIN ventas v ON v.id_vendedor = vd.id_vendedor
GROUP BY vd.id_vendedor, vd.nombre
HAVING AVG(v.unidades) > (SELECT AVG(unidades) FROM ventas)
ORDER BY promedio_unidades DESC, vd.nombre ASC;
\end{lstlisting}
\evidencia[0.7]{c09}{Vendedores cuyo promedio de unidades por venta supera el promedio general.}
\FloatBarrier

\subsection*{Consulta 10. Monto total vendido por cada vendedor}
\begin{lstlisting}
SELECT vd.nombre AS vendedor,
       SUM(v.unidades * b.precio) AS monto_total
FROM vendedores vd
LEFT JOIN ventas v ON v.id_vendedor = vd.id_vendedor
LEFT JOIN bicicletas b ON b.id_bicicleta = v.id_bicicleta
GROUP BY vd.id_vendedor, vd.nombre
ORDER BY monto_total DESC;
\end{lstlisting}
\textbf{Explicación solicitada:} Los vendedores sin ventas muestran \texttt{NULL} porque \texttt{SUM()} no dispone de valores no nulos para sumar. Para mostrar cero, se debe usar \texttt{COALESCE(SUM(v.unidades * b.precio), 0)}.

\evidencia[0.7]{c10}{Monto total vendido por cada vendedor, incluidos valores NULL.}
\FloatBarrier

\end{document}
